\section{Isolation and adjoint lifting}
\label{sec:isolation}

We now use a point-centered minimizing pair to construct an effective
rational SNC divisor whose coefficient equals the log discrepancy on a
nonempty reduced divisor over $x$, and is strictly smaller on every
other component meeting that divisor. This local configuration will
allow Kawamata--Viehweg vanishing to lift a nonzero value in the fiber
of $P=K_X+(n+1)L$ at $x$.

\subsection{Rational supporting data}

Suppose henceforth, for contradiction, that \(P\) is not generated
at \(x\).  Propositions~\ref{prop:minimizers} and
\ref{prop:point-center} give fixed \(t<n+1\) and \(m\) for which the
infimum \(f_{m,t}(x)\) is attained and every minimizing monomial pair
has center \(x\).  Choose such a pair \((s_1,\ldots,s_{N_m};v)\)
on a log resolution \(\pi:Y\to X\) of the divisors
\(D_i=\operatorname{div}(s_i)\).  Include all exceptional divisors in
the SNC divisor on \(Y\), and denote its components by \(E_j\).
Put
\[
 a_j=\A(E_j),\qquad
 z_j=\left(\frac{\ord_{E_j}(D_i)}{a_j}\right)_{i=1}^{N_m}.
\]
These vectors are rational and nonnegative, and \(a_j\) is a positive
integer because \(X\) is smooth.

For every component \(Z\) of an SNC stratum whose closure has image
containing \(x\), let \(J(Z)\) be its set of crossing components and
form the polytope
\begin{equation}
 C_Z=\conv\{z_j:j\in J(Z)\}.
 \label{eq:order-polytopes}
\end{equation}
Only nonempty sets \(J(Z)\) are used.  There are finitely many such
polytopes.  A convex representation \(z=\sum_j\lambda_jz_j\) is
realized by the monomial weights \(\lambda_j/a_j\); their discrepancy
is one and their section-order vector is \(z\).  They can be rescaled
to any positive discrepancy.  If some weights vanish, the center
enlarges to the appropriate stratum closure containing \(Z\).
Its image still contains \(x\).

Write
\[
 a=\A(v)>0,
 \qquad q=\left(\frac{v(s_i)}a\right)_{i=1}^{N_m}.
\]
The point \(q\) belongs to the union of the polytopes
\eqref{eq:order-polytopes}.  It belongs to none indexed by a stratum
with positive-dimensional image.  Indeed, a representation in such
a polytope, rescaled to discrepancy \(a\), would give the same
objective value and a center containing that positive-dimensional
image.  This would contradict Proposition~\ref{prop:point-center}.

Varying the scale of \(v\) differentiates the objective at its
minimum.  With
\begin{equation}
 b_i=\frac{t}{m}
       \frac{\exp(-v(s_i)/m)}{\sum_{h=1}^{N_m}\exp(-v(s_h)/m)},
 \label{eq:supporting-functional}
\end{equation}
this derivative gives
\begin{equation}
 b\cdot q=1,
 \qquad b_i>0,
 \qquad m\sum_i b_i=t<n+1.
 \label{eq:supporting-budget}
\end{equation}
For any \(C_Z\) containing \(q\), vary the normalized profile along
the segment from \(q\) to \(q'\in C_Z\), holding the discrepancy
equal to \(a\).  The one-sided derivative at the minimum is
nonnegative.  Since the objective decreases with each section order,
this says
\begin{equation}
 b\cdot q'\leq1\qquad(q'\in C_Z, q\in C_Z).
 \label{eq:supporting-inequality}
\end{equation}

The coefficients needed for vanishing must be rational. We must
approximate both $q$ and $b$ while preserving these supporting
inequalities and the exclusion of every polytope with
positive-dimensional image. The following elementary lemma preserves
exactly which indexed polytopes contain the profile; the family may
include repeated polytopes.

\begin{lemma}[Rational supporting data]
\label{lem:rational-support}
Let \(\mathcal C\) be a finite family of rational polytopes in
\(\R^N\), and let \(q\) belong to their union.  Suppose that a vector
\(b\in\R_{>0}^N\) satisfies
\[
 b\cdot q=1,
 \qquad b\cdot z\leq 1
 \quad\text{for every }z\in C\text{ whenever }q\in C\in\mathcal C.
\]
Fix an open neighborhood \(V\) of \(b\).  There are rational vectors
\(q^0\) arbitrarily close to \(q\) and
\(b^0\in V\cap\Q_{>0}^N\) such that
\begin{enumerate}
\item \(q^0\in C\) if and only if \(q\in C\), for each
      \(C\in\mathcal C\);
\item \(b^0\cdot q^0=1\), and \(b^0\cdot z\leq1\) on every
      \(C\in\mathcal C\) containing \(q^0\).
\end{enumerate}
\end{lemma}

\begin{proof}
For each polytope containing \(q\), let \(F_C\) be its smallest face
containing \(q\).  The point \(q\) lies in the relative interior of
\(F_C\).  Since the functional \(b\) attains its maximum at \(q\),
it has value one throughout \(F_C\).  The intersection
\[
 F=\bigcap_{C\ni q}F_C
\]
is a nonempty rational polytope, with \(q\) in its relative interior.
Its rational points are dense in it.  Choose a rational point \(q^0\)
of \(F\) sufficiently close to \(q\) to avoid every member of
\(\mathcal C\) not containing \(q\).  This is possible because those
members form a finite collection of closed sets disjoint from \(q\).
The incidence assertion follows, and \(b\cdot q^0=1\).

For this fixed rational \(q^0\), the conditions on a functional \(c\)
are the rational linear equation \(c\cdot q^0=1\) and the finitely
many rational linear inequalities \(c\cdot z\leq1\) at the vertices
of the polytopes containing \(q^0\).  Their feasible polyhedron
contains \(b\).  Rational points are dense in the smallest face of
this rational polyhedron containing \(b\).  A sufficiently close such
point belongs to \(V\) and has strictly positive coordinates; take it
as \(b^0\).
\end{proof}

Apply Lemma~\ref{lem:rational-support}, taking the neighborhood of
\(b\) small enough to preserve the strict budget in
\eqref{eq:supporting-budget}.  We obtain rational nonnegative
\(q^0\) and rational positive \(b^0\) with
\begin{equation}
 b^0\cdot q^0=1,
 \qquad m\sum_i b_i^0<n+1,
 \qquad b^0\cdot q'\leq1
 \quad(q'\in C_Z, q^0\in C_Z).
 \label{eq:rational-support-properties}
\end{equation}
The profile \(q^0\) belongs to no polytope with positive-dimensional
image.  In particular, \(q^0\ne0\).

A rational point in a rational convex hull has a rational convex
representation.  The corresponding rational monomial weights,
after clearing denominators, determine a primitive rational ray in
the SNC crossing.  Toroidal extraction of this ray gives a prime
divisor \(F_0\) whose normalized section-order vector is \(q^0\).
The monomial discrepancy formula gives equality in the retraction
inequality of Lemma~\ref{lem:retraction}.  Its center is \(x\), by
the exclusion of all positive-dimensional-image polytopes.  We may
retain this divisor on every subsequent common resolution: its
orders, discrepancy, and equality under retraction to \(Y\) depend
on the divisorial valuation, not on the higher model exhibiting it.

\subsection{An isolated discrepancy-equality divisor}

Let \(I=\{i:q_i^0>0\}\), and choose an integer \(M>0\) such that
\(M/q_i^0\) is an integer for every \(i\in I\).  Define the effective
ideal
\begin{equation}
 \mathfrak c=
 \sum_{i\in I}\cO_X\bigl(-(M/q_i^0)D_i\bigr).
 \label{eq:isolating-ideal}
\end{equation}
For any prime $F$, the order of this ideal divided by $M$ is the
minimum of the ratios $\ord_F(D_i)/q_i^0$, for $i\in I$.
At $F_0$ all these ratios equal $\A(F_0)$. A small contribution from
this minimum will preserve equality at $F_0$. In the coefficient
comparison below, equality in the perturbed bound will force these
ratios to agree; positivity of every $b_i^0$ will also control the
coordinates outside $I$.

Choose an integer \(u\) large enough that \(\mathfrak c(uL)\) is
globally generated.  Take a smooth projective resolution
\(\rho:T\to X\) dominating \(Y\), retaining \(F_0\), and
principalizing \(\mathfrak c\):
\[
 \mathfrak c\cO_T=\cO_T(-G).
\]
Arrange that \(G\), the pullbacks of the \(D_i\), and all exceptional
divisors have SNC support.  The divisor \(u\rho^*L-G\) is globally
generated, since it is the quotient of the pulled-back generating
system of \(\mathfrak c(uL)\).

For a sufficiently small rational number \(0<\eta<1\), put
\begin{equation}
 B=(1-\eta)\rho^*\left(\sum_i b_i^0D_i\right)
       +\frac{\eta}{M}G.
 \label{eq:perturbed-boundary}
\end{equation}
It is an effective rational divisor with SNC support.  Since
\(D_i\sim mL\), we have
\begin{align}
 (n+1)\rho^*L-B
 &\equiv
 \left(n+1-(1-\eta)m\sum_i b_i^0-\frac{\eta u}{M}\right)\rho^*L
 \notag\\
 &\hspace{2em}+\frac{\eta}{M}(u\rho^*L-G).
 \label{eq:boundary-nef-big}
\end{align}
The first coefficient is positive when \(\eta\) is sufficiently
small, by \eqref{eq:rational-support-properties}.  Thus this
difference is nef and big.

For a prime divisor \(F\) on \(T\), write
\[
 a_F=\A(F),\qquad
 r_F=\sum_j\ord_F(E_j)a_j\leq a_F.
\]
Here orders of divisors on \(Y\) mean orders of their pullbacks.
The number \(r_F\) is the discrepancy of the SNC retraction to
\(Y\).  Let \(S\) be the reduced sum of the primes in the chosen SNC
support satisfying all three conditions
\begin{equation}
 x\in\rho(F),\qquad
 \left(\frac{\ord_F(D_i)}{a_F}\right)_i=q^0,
 \qquad r_F=a_F.
 \label{eq:equality-divisor-definition}
\end{equation}
The retained divisor \(F_0\) is one of them, so \(S\ne0\).

\begin{lemma}
\label{lem:isolated-equality}
Every component of \(S\) maps to \(x\) and has coefficient \(a_F\)
in \(B\).  Every other prime in the SNC support meeting \(S\) has
coefficient strictly less than its log discrepancy.
\end{lemma}

\begin{proof}
If \(F\) is a component of \(S\), its retraction has normalized
profile \(q^0\), and the image of the retraction center contains
\(\rho(F)\), hence contains \(x\).  That image cannot be
positive-dimensional by our choice of \(q^0\).  Consequently
\(\rho(F)=\{x\}\).

Orders of a sum of ideals are minima of orders.  The definition of
\(\mathfrak c\) therefore gives, for every prime \(F\) on \(T\),
\begin{equation}
 \frac{\ord_F(G)}{M}
   =\min_{i\in I}\frac{\ord_F(D_i)}{q_i^0}.
 \label{eq:ideal-order-minimum}
\end{equation}
For a component of \(S\), this is \(a_F\).  Together with
\(b^0\cdot q^0=1\), Equation~\eqref{eq:perturbed-boundary} gives
\(\operatorname{coeff}_F B=a_F\).

Now let \(F\) meet a component \(F_S\) of \(S\), and choose a point
of intersection.  Its image \(y\) on \(Y\) lies over \(x\).
Take the stratum through \(y\) defined by all SNC components through
\(y\), and the corresponding polytope \(C\) from
\eqref{eq:order-polytopes}.  Every component with positive order
along either \(F\) or \(F_S\) contains the corresponding generic
center on \(Y\), and hence contains \(y\).  Both retraction profiles
therefore belong to \(C\), whenever their retractions are nonzero.
In particular \(q^0\in C\), so the supporting inequality in
\eqref{eq:rational-support-properties} applies on this common
polytope.

If \(r_F=0\), Lemma~\ref{lem:retraction} gives zero order on every
\(D_i\), and \eqref{eq:ideal-order-minimum} gives zero order on
\(G\).  Thus \(\operatorname{coeff}_F B=0<a_F\).  If \(r_F>0\),
write
\[
 q'=\left(\frac{\ord_F(D_i)}{r_F}\right)_i\in C,
 \qquad \mu=\min_{i\in I}\frac{q'_i}{q_i^0}.
\]
The positive numbers \(b_i^0q_i^0\), for \(i\in I\), sum to one.
Consequently
\begin{equation}
 \mu\leq\sum_{i\in I}b_i^0q'_i
       \leq b^0\cdot q'\leq1.
 \label{eq:rigid-minimum-inequality}
\end{equation}
Using \eqref{eq:ideal-order-minimum}, the exact coefficient ratio is
\begin{equation}
 \frac{\operatorname{coeff}_F B}{a_F}
   =\frac{r_F}{a_F}
       \bigl((1-\eta)b^0\cdot q'+\eta\mu\bigr)
   \leq1.
 \label{eq:boundary-discrepancy-ratio}
\end{equation}

Equality forces \(r_F=a_F\) and, since both coefficients in the
convex combination are positive, \(b^0\cdot q'=\mu=1\).
The equalities in \eqref{eq:rigid-minimum-inequality} then force
\(q'_i/q_i^0=1\) for every \(i\in I\).  Strict positivity of all
\(b_i^0\) also forces \(q'_i=0\) for \(i\notin I\).  Thus
\(q'=q^0\), including its zero coordinates.  Because \(F\) meets
\(F_S\) over \(x\), its image contains \(x\).  It now satisfies
all conditions in \eqref{eq:equality-divisor-definition} and is
itself a component of \(S\).  This proves the strict inequality for
every other component meeting \(S\).
\end{proof}

\subsection{Lifting and descent}\label{sec:final-lift}

The remaining argument follows the discrepancy-and-vanishing strategy
in \cite[\S1(1.6)]{fujita1993}, where strict discrepancy inequalities
along the divisors meeting an equality divisor allow a nonzero fiber
value to lift and descend. The construction above supplies that local
configuration on the reduced divisor $S$.

Define the integral Cartier divisor
\begin{equation}
 D=K_T+(n+1)\rho^*L-\lfloor B\rfloor.
 \label{eq:final-lifting-divisor}
\end{equation}
Its difference from the canonical divisor with fractional boundary is
\[
 D-(K_T+\{B\})=(n+1)\rho^*L-B.
\]
The fractional part \(\{B\}\) is an effective rational SNC boundary
with all coefficients in \([0,1)\), and the right side is nef and big
by \eqref{eq:boundary-nef-big}.  Theorem~\ref{thm:kv} therefore gives
\(H^1(T,\cO_T(D))=0\).  The exact sequence of the reduced Cartier
divisor \(S\) gives a surjection
\begin{equation}
 H^0(T,\cO_T(D+S))
       \longrightarrow H^0(S,\cO_S(D+S)).
 \label{eq:final-restriction-surjection}
\end{equation}

Compare the divisor upstairs with the desired adjoint bundle:
\begin{equation}
 C:=D+S-\rho^*P=K_{T/X}-\lfloor B\rfloor+S.
 \label{eq:final-correction}
\end{equation}
On a component \(F\) of \(S\), Lemma~\ref{lem:isolated-equality}
and the integrality of \(a_F\) give
\[
 \operatorname{coeff}_F C=(a_F-1)-a_F+1=0.
\]
For every other prime meeting \(S\), the same lemma gives
\(\lfloor\operatorname{coeff}_F B\rfloor\leq a_F-1\), and hence
\(\operatorname{coeff}_F C\geq0\).  Primes outside the chosen SNC
support contribute zero.  Thus \(C\) is effective on a neighborhood
of \(S\), with no component of \(S\) in its support.  Negative
components of \(C\) may occur elsewhere, but their supports are
disjoint from \(S\); removing those finitely many supports gives the
asserted neighborhood.

Since \(S\) is reduced and each of its components maps to \(x\),
the restriction \(\rho|_S\) factors scheme-theoretically through the
point \(x\).  Indeed, the pullback of a function vanishing at \(x\)
vanishes on every component of \(S\), and therefore vanishes on the
reduced scheme.  It follows that
\[
 \rho^*P|_S\simeq P|_x\otimes_{\C}\cO_S.
\]
Choose a nonzero element of the one-dimensional fiber \(P|_x\).
Multiply the corresponding constant section on \(S\) by the
canonical rational section of \(\cO_T(C)\).  This latter section is
regular near \(S\), because \(C\) is effective there, and is nonzero
at the generic point of each component of \(S\).  The product is
therefore a global section of \(\cO_S(D+S)\), nonzero at those generic
points.  Notice that no global effectiveness of \(C\) is required
for this restriction construction.

Lift this section using Equation~\eqref{eq:final-restriction-surjection}.
The positive part of \(C\) is exceptional over \(X\): on a
nonexceptional prime not belonging to \(S\), the
coefficient of \(C\) is \(-\lfloor\operatorname{coeff}_F B\rfloor
\leq0\), while on a component of \(S\) it is zero. Since \(S\) is
nonempty and every component maps to \(x\), any such component,
together with the imposed generic nonvanishing, satisfies the
hypotheses of Lemma~\ref{lem:descent}. The lift therefore descends to
a section \(s\in H^0(X,P)\) with \(s(x)\ne0\).

The argument also includes dimension one.  In that case a component
of \(S\) can be the nonexceptional point \(x\) itself; its
coefficient in \(C\) is still zero, and the same restriction and
descent argument applies.

We have contradicted the assumed failure of generation at \(x\).
Since a failure of global generation would occur at a closed point,
this proves Theorem~\ref{thm:main}.
